
\documentclass[aspectratio=169]{beamer}
\usetheme{Madrid}
\usecolortheme{wolverine}
\usepackage[utf8]{inputenc}
\usepackage{amsmath, amssymb}


\usepackage{tikz}

\title{GED Math Unit: Number Sense \& Problem Solving}
\author{Prepared for Instruction}
\date{\today}

\begin{document}

\begin{frame}
  \titlepage
\end{frame}



% =========================
\section{Lesson 1: Number Types and Place Value}

\begin{frame}{Objectives}
\begin{itemize}
  \item Identify and classify number types: whole numbers, integers, rationals, irrationals.
  \item Convert between fractions, decimals, and percentages.
  \item Use place value to round numbers appropriately.
\end{itemize}
\end{frame}

\begin{frame}{Key Ideas}
\begin{itemize}
  \item Place value: ones, tenths ($10^{-1}$), hundredths ($10^{-2}$), thousandths ($10^{-3}$), \dots
  \item Conversion: \(\text{percent} \to \text{decimal}\) by dividing by 100; \(\text{decimal} \to \text{fraction}\) by reading place value and simplifying.
  \item Number types: integers ($\dots,-2,-1,0,1,2,\dots$), rational (\(\frac{a}{b}, b\neq 0\)), irrational (non-repeating, non-terminating decimals).
\end{itemize}
\end{frame}

\begin{frame}{Worked Example 1}
\textbf{Convert } \(0.75\) \textbf{ to a fraction.}
\begin{align*}
0.75 &= \frac{75}{100} = \frac{3}{4}.
\end{align*}
\textbf{Answer: } \(\frac{3}{4}\).
\end{frame}

\begin{frame}{Worked Example 2}
\textbf{Simplify } \(\frac{18}{24}\).
\begin{align*}
\gcd(18,24) &= 6, \\
\frac{18}{24} &= \frac{18\div 6}{24\div 6} = \frac{3}{4}.
\end{align*}
\textbf{Answer: } \(\frac{3}{4}\).
\end{frame}

\begin{frame}{Worked Example 3}
\textbf{Round } \(0.03456\) \textbf{ to the nearest thousandth.}
\[
\text{Thousandth place} = 3\text{rd decimal digit} \to 0.034\ \text{(look at next digit 5, round up)} \Rightarrow 0.035.
\]
\textbf{Answer: } \(0.035\).
\end{frame}

\begin{frame}{Worked Example 4}
\textbf{Classify } \(-7,\, 3.14,\, \frac{1}{2}\).
\begin{itemize}
  \item \(-7\): integer (also rational).
  \item \(3.14\): terminating decimal \(\Rightarrow\) rational.
  \item \(\frac{1}{2}\): rational; as decimal \(0.5\).
\end{itemize}
\end{frame}

\begin{frame}{Worked Example 5}
\textbf{Convert } \(125\%\) \textbf{ to a decimal and fraction.}
\[
125\% = \frac{125}{100} = 1.25 = \frac{125}{100} = \frac{5}{4}.
\]
\textbf{Answer: } decimal \(1.25\), fraction \(\frac{5}{4}\).
\end{frame}

% =========================
\section{Lesson 2: Estimation and Mental Math}

\begin{frame}{Objectives}
\begin{itemize}
  \item Estimate sums, differences, products, and quotients.
  \item Use compatible numbers and front-end estimation.
  \item Judge reasonableness of answers.
\end{itemize}
\end{frame}

\begin{frame}{Worked Example 1}
\textbf{Estimate } \(234 + 567\).
\[
234 \approx 200,\quad 567 \approx 600 \Rightarrow 200 + 600 = 800.
\]
\textbf{Reasonable exact sum } \(= 801\), so estimate \(800\) is close.
\end{frame}

\begin{frame}{Worked Example 2}
\textbf{Estimate } \(82 \times 49\) \textbf{ using compatible numbers.}
\[
82 \approx 80,\quad 49 \approx 50 \Rightarrow 80 \times 50 = 4000.
\]
\textbf{Reasonableness: } exact \(82\times49=4018\).
\end{frame}

\begin{frame}{Worked Example 3}
\textbf{Approximate } \(1278 \div 6\).
\[
1278 \approx 1200,\quad 1200\div 6 = 200.
\]
Exact \(=213\), so estimate \(200\) is reasonable.
\end{frame}

\begin{frame}{Worked Example 4}
\textbf{Is } \(49\times 21 \approx 1000\) \textbf{ reasonable?}
\[
50\times 20 = 1000 \quad\text{(close)},\quad \text{exact }=1029.\ \text{Yes, reasonable.}
\]
\end{frame}

\begin{frame}{Worked Example 5}
\textbf{Round } \(23.746\) \textbf{ to nearest tenth and whole.}
\[
\text{Tenth: look at hundredths (4)} \Rightarrow 23.7;\quad \text{Whole: look at tenths (7)} \Rightarrow 24.
\]
\end{frame}

% =========================
\section{Lesson 3: Ratios, Proportions, Percents}

\begin{frame}{Objectives}
\begin{itemize}
  \item Write and simplify ratios.
  \item Solve proportions by cross-multiplication.
  \item Compute percent, percent change, and discounts.
\end{itemize}
\end{frame}

\begin{frame}{Worked Example 1}
\textbf{Write and simplify the ratio of 6 apples to 9 oranges.}
\[
6:9 = \frac{6}{9} = \frac{2}{3} \Rightarrow 2:3.
\]
\end{frame}

\begin{frame}{Worked Example 2}
\textbf{Solve } \(\frac{3}{4} = \frac{x}{16}\).
\[
3\cdot 16 = 4x \Rightarrow 48 = 4x \Rightarrow x=12.
\]
\end{frame}

\begin{frame}{Worked Example 3}
\textbf{Find } \(25\%\) \textbf{ of } \(120\).
\[
25\% = \frac{1}{4} \Rightarrow \frac{1}{4}\cdot 120 = 30.
\]
\end{frame}

\begin{frame}{Worked Example 4}
\textbf{Price increases from \$40 to \$50. Find percent increase.}
\[
\text{Change} = 50-40=10,\quad \frac{10}{40}=0.25=25\%.
\]
\end{frame}

\begin{frame}{Worked Example 5}
\textbf{\$30 shirt on sale for 20\% off. Find sale price.}
\[
\text{Discount} = 0.20\times 30 = 6,\quad \text{Sale price}=30-6=\$24.
\]
\end{frame}

% =========================
\section{Lesson 4: Problem-Solving Strategies}

\begin{frame}{Objectives}
\begin{itemize}
  \item Apply a four-step process: Understand, Plan, Solve, Check.
  \item Translate word problems into equations.
  \item Evaluate the reasonableness of solutions.
\end{itemize}
\end{frame}

\begin{frame}{Worked Example 1}
\textbf{``Twice a number is 10 more than 8.'' Find the number.}
\[
2x = 8 + 10 = 18 \Rightarrow x=9.
\]
\textbf{Check: } \(2\cdot 9=18\) matches.
\end{frame}

\begin{frame}{Worked Example 2}
\textbf{Recipe uses 3 cups flour for 12 cookies. How much for 20?}
\[
\text{Rate} = \frac{3}{12}=\frac{1}{4}\ \text{cup per cookie} \Rightarrow 20\times \frac{1}{4} = 5\ \text{cups}.
\]
\end{frame}

\begin{frame}{Worked Example 3}
\textbf{5 pencils cost \$2.25. Cost for 12 pencils?}
\[
\text{Unit price} = \frac{2.25}{5} = \$0.45,\quad 12\times 0.45 = \$5.40.
\]
\end{frame}

\begin{frame}{Worked Example 4}
\textbf{Earn \$15/hr for 35 hours. How much after 15\% tax?}
\[
\text{Gross} = 15\times 35 = \$525,\quad \text{Tax} = 0.15\times 525 = \$78.75,
\]
\[
\text{Net} = 525 - 78.75 = \$446.25.
\]
\end{frame}

\begin{frame}{Worked Example 5}
\textbf{Car travels 240 miles in 4 hours. Find speed. Then test 250 miles in 4.5 hrs.}
\[
\text{Speed} = \frac{240}{4}=60\ \text{mph}.
\]
\[
\text{At }60\ \text{mph for }4.5\text{ h} \Rightarrow 60\times 4.5 = 270\ \text{mi}. \ \ 250\text{ mi is reasonable but a bit less.}
\]
\end{frame}

% =========================
\section{Lesson 5: Real-World Word Problems}

\begin{frame}{Objectives}
\begin{itemize}
  \item Model mixture, distance-rate-time, work, and money problems.
  \item Interpret real-world scenarios and select operations.
\end{itemize}
\end{frame}

\begin{frame}{Worked Example 1}
\textbf{Distance: } A car travels 60 mph for 3 hours. How far?
\[
d=rt=60\times 3 = 180\ \text{miles}.
\]
\end{frame}

\begin{frame}{Worked Example 2}
\textbf{Tank holds 50 L and is 40\% full. How many liters?}
\[
0.40 \times 50 = 20\ \text{L}.
\]
\end{frame}

\begin{frame}{Worked Example 3}
\textbf{Two workers finish a job in 6 h. One alone takes 10 h. How long for the other alone?}
Let rates be \(A=\frac{1}{10}\) job/h and \(B=\frac{1}{x}\) job/h.
\[
A+B=\frac{1}{6} \Rightarrow \frac{1}{10} + \frac{1}{x} = \frac{1}{6}
\Rightarrow \frac{1}{x} = \frac{1}{6}-\frac{1}{10}=\frac{5-3}{30}=\frac{2}{30}=\frac{1}{15}.
\]
Thus \(x=15\) hours.
\end{frame}

\begin{frame}{Worked Example 4}
\textbf{8\% sales tax on a \$45 item. What is the total?}
\[
\text{Tax} = 0.08\times 45 = \$3.60,\quad \text{Total} = 45 + 3.60 = \$48.60.
\]
\end{frame}

\begin{frame}{Worked Example 5}
\textbf{Reading a bar graph of monthly spending. Identify highest category.}
\begin{itemize}
  \item Strategy: Compare bar heights (or values) and select the maximum.
  \item Always read axes labels and units before comparing.
\end{itemize}
\textit{(Insert class-specific data graphic if available.)}
\end{frame}

% =========================
\section{Practice \& Reflection}

\begin{frame}{Practice Prompts (Optional Add-on)}
\begin{enumerate}
  \item Convert \(0.375\) to a fraction in simplest form.
  \item Estimate \(389\times 21\) using compatible numbers.
  \item Solve \(\frac{7}{8} = \frac{y}{24}\).
  \item A jacket listed at \$64 has 30\% off and 8\% tax. Find the total price.
  \item A train travels \(45\) mph. How long to go \(157.5\) miles?
\end{enumerate}
\end{frame}

\begin{frame}{Exit Ticket}
\begin{itemize}
  \item One concept I mastered today is \underline{\hspace{4cm}}.
  \item One question I still have is \underline{\hspace{4cm}}.
  \item My strategy for solving word problems is \underline{\hspace{4cm}}.
\end{itemize}
\end{frame}

\end{document}
